\documentclass[11pt]{amsart}
\usepackage{amsmath,amssymb,amsthm}
\usepackage[margin=1.1in]{geometry}
\usepackage{array,longtable,booktabs}
\usepackage{hyperref}

\newtheorem{theorem}{Theorem}
\newtheorem{lemma}[theorem]{Lemma}
\newtheorem{proposition}[theorem]{Proposition}
\theoremstyle{remark}
\newtheorem{remark}[theorem]{Remark}

\newcommand{\Inf}{\mathrm{Inf}}
\newcommand{\Var}{\mathrm{Var}}
\newcommand{\E}{\mathbb{E}}

\title{Four equivalent forms of the Fourier entropy-influence conjecture}
\author{Yitzchak Shmalo}
\address{Einstein Institute of Mathematics, The Hebrew University of Jerusalem, Jerusalem, Israel}
\date{}

\begin{document}

\begin{abstract}
Let $f$ be a Boolean function with spectral distribution $\mu(S)=\hat f(S)^2$, and let $\gamma(T)=\sum_{U\supseteq
T}\mu(U)$ be the generalized influence of a set $T$. The Fourier entropy splits as $H(f)=A+B$ with
$A=\E_\mu\log_2(1/\gamma(S))$ and $B=\E_\mu\log_2(\gamma(S)/\mu(S))\ge 0$. Since derivatives of Boolean functions take
values in a lattice, $A\le H(f)/2+I(f)-\Var(f)$, and so the Fourier entropy-influence conjecture is equivalent to
$B\le C\,I(f)$. We give three further equivalent forms: a bound on the cross-entropy of the spectral distribution against
the noisy generalized influences $\|T_{\sqrt t}\,\partial_T f\|_2^2$ at a fixed $t<1$; a bound on the integral over
$t\in[\tfrac12,1]$ of the expected number of elements a noise-tilted posterior adds to the spectral sample; and a bound
of $O(1)$ bits on the log-loss of predicting a uniformly chosen element of the size-biased spectral sample from the
others, which also has a variational form over all predictors. Form (c) gives the unconditional bound
$H(f)\le 4I(f)-2\Var(f)+2M(f)$, where $M(f)$ is the expected number of such added elements at $t=1$.
\end{abstract}

\subjclass[2020]{06E30, 42C10, 94A17}
\keywords{Fourier entropy-influence conjecture, Boolean functions, Fourier entropy, total influence}

\maketitle

\section{Introduction}\label{sec:intro}

Let $f:\{-1,1\}^n\to\{-1,1\}$ have Fourier expansion $f=\sum_S\hat f(S)\chi_S$. By Parseval $\mu(S)=\hat f(S)^2$ is a
probability distribution on subsets of $[n]$, and we write $U$ for a random set with law $\mu$. Put
\[
H(f)=\sum_S\mu(S)\log_2\frac1{\mu(S)},\quad \Inf_i(f)=\sum_{S\ni i}\mu(S),\quad I(f)=\sum_S|S|\mu(S),\quad
\Var(f)=1-\mu(\emptyset).
\]
The Fourier entropy-influence (FEI) conjecture of Friedgut and Kalai~\cite{FK96} asks for a universal $C$ with
$H(f)\le C\,I(f)$ for all Boolean $f$. It is known for symmetric functions and read-once decision
trees~\cite{OWZ11} and is preserved under suitable compositions~\cite{OT13}. There are bounds in terms of
decision-tree and communication complexity measures~\cite{CKLS,WWW}, and it holds for linear threshold
functions~\cite{Shm2}. For every Boolean $f$, Han~\cite{Han} proved $H(f)\le(2+3/\ln 2)\,I(f)+\sum_i\Inf_i(f)\log_2(1/\Inf_i(f))$,
and Li and Han~\cite{LH25} lowered the constant in front of $I(f)$ to the sharp value $1/\ln 2$.
In~\cite{Shm1} the conjecture was shown to be equivalent to a one-coordinate inequality, and to the bound
$\sum_iE_i(f)\le C\,I(f)$ on the erasure entropies (in the sense of Verd\'u and Weissman~\cite{VW08})
\[
E_i(f)=H\bigl(\mathbf 1[i\in U]\bigm| U\setminus\{i\}\bigr).
\]
The one-coordinate route goes back to Han~\cite[Remarks~3.1--3.2]{Han}, and the splitting of the entropy drop at one
coordinate into $E_i$ and a Jensen--Shannon term is equation~(32) of Gonz\'alez, MacManus and Pereyra~\cite{GMP}.
Here we give four more equivalent forms. All of them use the derivatives
\[
\partial_T f=\sum_{U\supseteq T}\hat f(U)\chi_{U\setminus T},\qquad
\gamma_t(T)=\|T_{\sqrt t}\,\partial_T f\|_2^2=\sum_{U\supseteq T}\mu(U)\,t^{|U\setminus T|}\quad(0\le t\le 1),
\]
where $T_\rho$ is the noise operator. For $t=1$ we write $\gamma(T)=\gamma_1(T)=\|\partial_Tf\|_2^2$; this is the
generalized influence $I_T[f]$ of Kelman, Kindler, Lifshitz, Minzer and Safra~\cite{KKLMS}. Note $\gamma_0=\mu$,
$\gamma(\emptyset)=1$ and $\mu\le\gamma_t\le\gamma\le 1$. Sums over $S$ below run over the support of $\mu$. Put
\[
A=\sum_S\mu(S)\log_2\frac1{\gamma(S)},\qquad B=\sum_S\mu(S)\log_2\frac{\gamma(S)}{\mu(S)},\qquad H(f)=A+B,
\]
and for $0<t<1$ let $A_t=\sum_S\mu(S)\log_2(1/\gamma_t(S))$. Let
\[
m_t(T)=\frac{\sum_{U\supseteq T}|U\setminus T|\,\mu(U)\,t^{|U\setminus T|}}{\gamma_t(T)},\qquad
M_t(f)=\sum_S\mu(S)\,m_t(S),
\]
so that $m_t(T)$ is the mean of $|U\setminus T|$ when $U\supseteq T$ is drawn with probability proportional to
$\mu(U)t^{|U\setminus T|}$, and write $M(f)=M_1(f)$. Finally let $Q$ be the law on pairs $(T,i)$ with $i\notin T$ given
by $Q(T,i)=\mu(T\cup\{i\})/I(f)$; it is the law of $(U\setminus\{i\},i)$ when $U$ is drawn size-biased and $i$ is a
uniform element of $U$. Put
\[
L(f)=\E_{(T,i)\sim Q}\log_2\frac{\mu(T)+\mu(T\cup\{i\})}{\mu(T\cup\{i\})},
\]
the log-loss in bits of predicting that $i\in U$ from $U\setminus\{i\}$.

\begin{theorem}\label{thm:main}
For every Boolean function $f$ and every $0<t<1$,
\begin{enumerate}
\item[(a)] $H(f)\le 2B+2I(f)-2\Var(f)$ and $0\le B\le H(f)$;
\item[(b)] $A_t\le H(f)\le A_t+I(f)\log_2\frac1{1-t}$;
\item[(c)] $B=\frac1{\ln 2}\int_0^1M_s(f)\,\frac{ds}{s}$, the function $s\mapsto m_s(T)$ is nondecreasing for every
$T$, and
\[
H(f)\le 4I(f)-2\Var(f)+\frac{2}{\ln2}\int_{1/2}^1M_s(f)\,\frac{ds}{s}\le 4I(f)-2\Var(f)+2M(f);
\]
\item[(d)] $I(f)\,L(f)\le\sum_iE_i(f)\le I(f)\,L(f)+I(f)/\ln 2$.
\end{enumerate}
Consequently, with~\cite[Corollary~5]{Shm1} for the last item, the following are equivalent: the FEI conjecture; $B\le C\,I(f)$ for all $f$; $A_t\le C\,I(f)$ for all $f$
(for one, or every, fixed $t\in(0,1)$); $\int_{1/2}^1M_s(f)\,ds\le C\,I(f)$ for all $f$; and $L(f)\le C$ for all $f$.
\end{theorem}

Each form isolates a different piece of the difficulty. Form (a) says that the part of the entropy carried by the size of
the generalized influences is harmless, and what must be bounded is the relative entropy of $\mu(S)$ against
$\gamma(S)$, which compares a coefficient with the total weight above it. Form (b) replaces the exact coefficients by
their noisy upper sets, at any fixed noise level. Form (c) writes $B$ as the expected number of elements a noise-tilted
posterior adds to the spectral sample, and shows that only noise levels near $t=1$ matter. Form (d) reduces the
conjecture to the prediction of a single element of the spectral sample.

The equivalences hold for the class of all Boolean functions, not function by function. For instance, when $\hat f$ is
supported on sets of odd size, as for the majority function, the parity of $|U\setminus\{i\}|$ determines whether
$i\in U$, so $E_i(f)=0$ for every $i$ and $L(f)=0$, while $H(f)$ may be large compared with $I(f)$.

\section{The split}\label{sec:split}

\begin{lemma}\label{lem:gran}
For every $T\ne\emptyset$, $\mu(T)\le 4^{|T|-1}\gamma(T)^2$.
\end{lemma}

\begin{proof}
This is the usual granularity argument for derivatives of Boolean functions~\cite[Chapter~2]{OD14}. For
$x\in\{-1,1\}^n$, $\partial_Tf(x)=2^{-|T|}\sum_{y\in\{-1,1\}^T}\bigl(\prod_{j\in T}y_j\bigr)f(y,x_{[n]\setminus T})$, a sum of
$2^{|T|}$ terms $\pm1$ divided by $2^{|T|}$. The sum is even, so $\partial_Tf$ takes values in $2^{1-|T|}\mathbb Z$ and
$(\partial_Tf)^2\ge 2^{1-|T|}|\partial_Tf|$ pointwise. Hence
$\gamma(T)=\E(\partial_Tf)^2\ge 2^{1-|T|}\E|\partial_Tf|\ge 2^{1-|T|}|\hat f(T)|$.
\end{proof}

\begin{proof}[Proof of Theorem~\ref{thm:main}(a)]
Since $\gamma\le1$ and $\gamma\ge\mu$, $A\ge0$ and $B\ge0$, so $B\le H(f)$. By Lemma~\ref{lem:gran}, for $S\ne\emptyset$
in the support, $\log_2(1/\gamma(S))\le |S|-1+\tfrac12\log_2(1/\mu(S))$, while $S=\emptyset$ contributes $0$ to $A$. So
\[
A\le\sum_{S\ne\emptyset}\mu(S)(|S|-1)+\tfrac12H(f)=I(f)-\Var(f)+\tfrac12H(f),
\]
and $H(f)=A+B$ gives the claim.
\end{proof}

The constant in Lemma~\ref{lem:gran} is attained: for a dictator $f=\chi_{\{1\}}$ and $T=\{1\}$ both sides equal $1$.

\section{Noisy upper sets}\label{sec:thin}

\begin{proof}[Proof of Theorem~\ref{thm:main}(b)]
Let $\nu_t(T)=(1-t)^{|T|}\gamma_t(T)=\sum_{U\supseteq T}\mu(U)(1-t)^{|T|}t^{|U\setminus T|}$. This is the law of the
random set obtained from $U$ by keeping each element independently with probability $1-t$, so $\sum_T\nu_t(T)=1$.
Gibbs' inequality gives
\[
H(f)\le\sum_S\mu(S)\log_2\frac1{\nu_t(S)}=A_t+I(f)\log_2\frac1{1-t}.
\]
The other bound is $\gamma_t\ge\mu$.
\end{proof}

At $t=1$ only the lower bound survives: $A\le H(f)$, and indeed $A\le H(f)/2+I(f)-\Var(f)$ by part (a). So the bound $A\le
C\,I(f)$ is a consequence of FEI but is not known to imply it.

\section{The tilted excess}\label{sec:tilt}

\begin{proof}[Proof of Theorem~\ref{thm:main}(c)]
Fix $T$ in the support of $\mu$. Then $\frac{d}{ds}\ln\gamma_s(T)=m_s(T)/s$ for $0<s\le1$, and $m_s(T)=O(s)$ as $s\to0$.
Differentiating once more, $s\frac{d}{ds}m_s(T)$ is the variance of $|U\setminus T|$ under the tilted law in the
definition of $m_s$, so $m_s(T)$ is nondecreasing. Since $\gamma_0(T)=\mu(T)$,
\[
\log_2\frac{\gamma(T)}{\mu(T)}=\frac1{\ln2}\int_0^1m_s(T)\,\frac{ds}{s},\qquad
\log_2\frac{1}{\gamma_\lambda(T)}=\log_2\frac1{\gamma(T)}+\frac1{\ln2}\int_\lambda^1m_s(T)\,\frac{ds}{s}.
\]
Averaging the first identity over $\mu$ (the integrand is nonnegative) gives the formula for $B$. Averaging the second
gives $A_\lambda=A+\frac1{\ln2}\int_\lambda^1M_s\,ds/s$. Part (b) at $t=\lambda$ and the bound on $A$ from the proof of
part (a) give
\[
H(f)\le\tfrac12H(f)+I(f)-\Var(f)+\frac1{\ln2}\int_\lambda^1M_s(f)\,\frac{ds}{s}+I(f)\log_2\frac1{1-\lambda}.
\]
Taking $\lambda=\tfrac12$ gives the first inequality, and $M_s\le M_1$ gives the second.
\end{proof}

For the equivalence in the theorem, one direction is the inequality just proved, with $ds/s\le 2\,ds$ on $[\tfrac12,1]$.
For the other, $\int_{1/2}^1M_s\,ds\le\int_0^1M_s\,ds/s=B\ln2\le H(f)\ln 2$.

The quantity $M(f)$ is the expected number of elements that a draw from $\mu$ conditioned to contain $S$ adds to $S$,
when $S$ is itself drawn from $\mu$. The bound $H(f)\le 4I(f)-2\Var(f)+2M(f)$ holds for every $f$; the bound
$M(f)\le C\,I(f)$ for all Boolean $f$ would imply the conjecture, and it is open whether it holds.

\section{Predicting one element}\label{sec:pred}

\begin{proof}[Proof of Theorem~\ref{thm:main}(d)]
Fix $i$, and for $T\not\ni i$ write $p=\mu(T)$ and $q=\mu(T\cup\{i\})$. The contribution of the pair to $E_i(f)$ is
\[
p\log_2\frac{p+q}{p}+q\log_2\frac{p+q}{q}.
\]
Summing the second term over $i$ and $T$ gives $I(f)L(f)$. The first term lies in $[0,q/\ln 2]$, and $\sum_i\sum_T
q=I(f)$.
\end{proof}

Combined with~\cite[Corollary~5]{Shm1}, which says that FEI with constant $C$ implies $\sum_iE_i\le C\,I(f)$ and that
$\sum_iE_i\le C\,I(f)$ for all $f$ implies FEI with constant $C+2/\ln2$, part (d) shows that FEI holds if and only if
$\sup_fL(f)<\infty$; precisely, $L\le C$ for all $f$ gives $H(f)\le(C+3/\ln2)I(f)$.

The loss $L$ has a variational form. For any function $p_i(T)\ge0$ on pairs, put
\[
K_p=\frac1{I(f)}\sum_{i}\sum_{T\not\ni i}\bigl(\mu(T)+\mu(T\cup\{i\})\bigr)\,p_i(T).
\]

\begin{proposition}\label{prop:var}
For every $p$, $L(f)\le\log_2K_p+\E_Q\log_2\frac1{p_i(T)}$, with equality when
$p_i(T)=\mu(T\cup\{i\})/(\mu(T)+\mu(T\cup\{i\}))$.
\end{proposition}

\begin{proof}
$R(T,i)=(\mu(T)+\mu(T\cup\{i\}))\,p_i(T)/(K_pI(f))$ is a probability law on pairs, and Gibbs' inequality
$\E_Q\log_2(Q/R)\ge0$ is the claim. At the stated $p$, $R=Q$ and $K_p=1$.
\end{proof}

So a predictor that guesses ``$i\in U$'' with probability $p_i(T)$, and whose normalized total guess $K_p$ is bounded,
bounds $L$. The constant predictor $p\equiv\tfrac12$ has $K_p=n/(2I(f))$ and gives $L(f)\le\log_2(n/I(f))$, hence
$\sum_iE_i(f)\le I(f)\bigl(\log_2(n/I(f))+1/\ln2\bigr)$. A bound uniform in $n$ needs a predictor that uses $T$.

\begin{remark}\label{rem:centre}
Multiplying $f$ by a character $\chi_R$ moves the spectrum, $\widehat{\chi_Rf}(S)=\hat f(S\triangle R)$. It maps the
pairs $\{T,T\cup\{i\}\}$ to pairs, so every $E_i$ is unchanged, while $\Inf_i$ becomes $1-\Inf_i$ for $i\in R$. Choosing
$R=\{i:\Inf_i(f)>\tfrac12\}$ shows that the bound $\sum_iE_i(f)\le C\,I(f)$ for all $f$ is equivalent to
$\sum_iE_i(f)\le C\sum_i\min(\Inf_i(f),1-\Inf_i(f))$ for all $f$; with~\cite[Corollary~5]{Shm1}, so is FEI. (For $H$
itself the same move is Lemma~2.2 of~\cite{Shm2}.) Since $D_if$ takes values in $\{-1,0,1\}$, $\Inf_i(f)=\Pr[D_if\ne0]$
and $\hat f(\{i\})=\mathbf E\,D_if$, so $\hat f(\{i\})^2\le\Inf_i(f)^2$ and
with $W_1(f)=\sum_i\hat f(\{i\})^2$,
\[
I(f)-W_1(f)\ge\sum_i\Inf_i(f)\bigl(1-\Inf_i(f)\bigr)\ge\tfrac12\sum_i\min\bigl(\Inf_i(f),1-\Inf_i(f)\bigr).
\]
Hence FEI is also equivalent to $\sum_iE_i(f)\le C\,(I(f)-W_1(f))$ for all $f$, with $C$ changed by a factor at most $2$.
\end{remark}

\section{Checks}\label{sec:ex}

All identities and inequalities of Theorem~\ref{thm:main}, Lemma~\ref{lem:gran}, Proposition~\ref{prop:var} for the
constant predictor and the invariance in Remark~\ref{rem:centre} were checked with exact Walsh--Hadamard transforms on
$119$ functions: random functions on $3$ to $9$ variables with biases $\tfrac12$, $0.15$ and $0.03$, the AND of eight
variables, the tribes functions of widths $2$ and $3$ on eight and nine variables, and majority on nine variables. The
ratio in Lemma~\ref{lem:gran} has maximum exactly $1$; the identities $H=A+B$ and the invariance of $\sum_iE_i$ hold to
$2\cdot10^{-15}$ and the integral formula for $B$ to $8\cdot10^{-7}$ with $80$-point Gauss--Legendre quadrature; the
inequalities of parts (a) and (b) (at $t=0.25,0.5,0.75,0.9$) and the bound with $M(f)$ hold on all of these functions,
with smallest slack $0.022$, $0.0055$ and $0.16$ respectively (a character has slack $0$ in (a) and (b)); and $m_s$ is nondecreasing on the quadrature grid for every function. For majority on nine
variables, $\sum_iE_i=L=0$ while $H/I=2.137$.

\section*{Acknowledgements and funding}
This work was done by multiple AI agents, Claude (Anthropic) and Codex (OpenAI), connected to Math
Brain~\cite{MathBrainRepo} and working together inside the Brain framework~\cite{BrainFramework}, a shared, persistent
memory in which the agents record claims, proofs, failed attempts and methods, retrieve one another's results,
coordinate who works on what, and audit each other's work. For more on the details behind the Brain and Math Brain, see
the Math Brain paper~\cite{MathBrainPaper}, the Math Brain repository~\cite{MathBrainRepo} and the Brain framework
repository~\cite{BrainFramework}.

The research presented in this paper was supported by the European Research Council (ERC) under the European Union's
Horizon Europe research and innovation programme (grant agreement No.~101041711), by the Simons Foundation as part of
the Collaboration on the Mathematical and Scientific Foundations of Deep Learning, by Heights Labs, by Convex Nexus
Capital, by the Israel Science Foundation (grant number 2258/19), and by the Israel Science Foundation (ISF Grant
4101/25). Computational resources were provided in part by Google Cloud through its education credits programme (credit
530550103).

\section*{Declaration of competing interest}
The author holds equity in Heights Labs and in Convex Nexus Capital and has a management role at Convex Nexus Capital.
These entities are included in the funding declaration above.

\section*{Declaration of generative AI and AI-assisted technologies}
The paper comes out of Math Brain, a project on the use of AI systems in research mathematics~\cite{BrainFramework,MathBrainRepo}. In November 2026 a
Claude Fable agent (Anthropic), connected to the Brain and working alongside Codex agents running OpenAI's GPT-6.1 on
the same problem, found the split of part (a), the noisy form of part (b), the integral formula of part (c) and the
predictor form of part (d). Claude Opus 5.5 (Anthropic) searched the literature, rederived each proof, found the bound
with $M(f)$, wrote the numerical check, prepared the manuscript and audited it. The author directed the work and is responsible for the content of the paper.

\section*{Code and data availability}
Code and outputs: \url{https://github.com/yspennstate/fei-equivalent-forms}.

The paper is part of the Math Brain project, whose public repository is
\url{https://github.com/yspennstate/math-brain-public}~\cite{MathBrainRepo}.

\appendix
\section{Notation}\label{app:notation}

Symbols are listed in the order in which they first appear, and the last column gives that place. The letter $T$
denotes a set and $T_\rho$ the noise operator. The letters $p$ and $R$ are reused; for these the last column gives the
place where each meaning first appears.

\begin{longtable}{@{}>{\raggedright\arraybackslash}p{0.19\textwidth}>{\raggedright\arraybackslash}p{0.54\textwidth}>{\raggedright\arraybackslash}p{0.2\textwidth}@{}}
\toprule
Symbol & Meaning & First appears \\
\midrule
\endhead
$f$ & a Boolean function $\{-1,1\}^n\to\{-1,1\}$ & Section~\ref{sec:intro} \\
$\hat f(S)$, $\chi_S$ & the Fourier coefficient of $f$ at $S\subseteq[n]$, and the character $\prod_{i\in S}x_i$ & Section~\ref{sec:intro} \\
$\mu$ & the spectral distribution, $\mu(S)=\hat f(S)^2$ & Section~\ref{sec:intro} \\
$U$ & a random set with law $\mu$ (the spectral sample) & Section~\ref{sec:intro} \\
$H(f)$ & the Fourier entropy $\sum_S\mu(S)\log_2(1/\mu(S))$ & Section~\ref{sec:intro} \\
$\Inf_i(f)$ & the influence of $i$, $\sum_{S\ni i}\mu(S)$ & Section~\ref{sec:intro} \\
$I(f)$ & the total influence $\sum_S|S|\mu(S)$ & Section~\ref{sec:intro} \\
$\Var(f)$ & the variance $1-\mu(\emptyset)$ & Section~\ref{sec:intro} \\
$C$ & a universal constant, as in $H(f)\le C\,I(f)$ & Section~\ref{sec:intro} \\
$E_i(f)$ & the erasure entropy $H(\mathbf 1[i\in U]\mid U\setminus\{i\})$ & Section~\ref{sec:intro} \\
$\partial_Tf$ & the derivative $\sum_{U\supseteq T}\hat f(U)\chi_{U\setminus T}$ & Section~\ref{sec:intro} \\
$T_\rho$ & the noise operator, $T_\rho\chi_S=\rho^{|S|}\chi_S$ & Section~\ref{sec:intro} \\
$t$ & a noise parameter in $[0,1]$ & Section~\ref{sec:intro} \\
$\gamma_t(T)$ & $\|T_{\sqrt t}\,\partial_Tf\|_2^2=\sum_{U\supseteq T}\mu(U)\,t^{|U\setminus T|}$ & Section~\ref{sec:intro} \\
$\gamma(T)$ & $\gamma_1(T)=\|\partial_Tf\|_2^2$, the generalized influence of $T$ & Section~\ref{sec:intro} \\
$A$, $B$ & $\sum_S\mu(S)\log_2(1/\gamma(S))$ and $\sum_S\mu(S)\log_2(\gamma(S)/\mu(S))$, so that $H(f)=A+B$ & Section~\ref{sec:intro} \\
$A_t$ & $\sum_S\mu(S)\log_2(1/\gamma_t(S))$ & Section~\ref{sec:intro} \\
$m_t(T)$ & the mean of $|U\setminus T|$ when $U\supseteq T$ is drawn with probability proportional to $\mu(U)t^{|U\setminus T|}$ & Section~\ref{sec:intro} \\
$M_t(f)$, $M(f)$ & $\sum_S\mu(S)\,m_t(S)$, and $M_1(f)$ & Section~\ref{sec:intro} \\
$Q$ & the law $Q(T,i)=\mu(T\cup\{i\})/I(f)$ on pairs $(T,i)$ with $i\notin T$ & Section~\ref{sec:intro} \\
$L(f)$ & the log-loss $\E_Q\log_2\bigl((\mu(T)+\mu(T\cup\{i\}))/\mu(T\cup\{i\})\bigr)$ of predicting that $i\in U$ from $U\setminus\{i\}$ & Section~\ref{sec:intro} \\
$\nu_t(T)$ & $(1-t)^{|T|}\gamma_t(T)$, the law of the set that keeps each element of $U$ with probability $1-t$ & Section~\ref{sec:thin} \\
$s$ & the integration variable in part (c) & Theorem~\ref{thm:main} \\
$\lambda$ & the noise level at which part (b) is applied in the proof of part (c) & Section~\ref{sec:tilt} \\
$p$, $q$ & $\mu(T)$ and $\mu(T\cup\{i\})$ for a fixed $i$ and $T\not\ni i$, in the proof of part (d) & Section~\ref{sec:pred} \\
$p_i(T)$, $K_p$ & a predictor, a nonnegative function on pairs, and its normalized total $\frac1{I(f)}\sum_i\sum_{T\not\ni i}(\mu(T)+\mu(T\cup\{i\}))\,p_i(T)$ & Section~\ref{sec:pred} \\
$R(T,i)$, $R$ & the probability law $(\mu(T)+\mu(T\cup\{i\}))\,p_i(T)/(K_pI(f))$ on pairs, in the proof of Proposition~\ref{prop:var}; in Remark~\ref{rem:centre}, a set, with $\chi_Rf$ the shifted function, and then $\{i:\Inf_i(f)>\tfrac12\}$ & Proposition~\ref{prop:var}; Remark~\ref{rem:centre} \\
$D_if$ & $\tfrac12(f_{i\to1}-f_{i\to-1})$, with values in $\{-1,0,1\}$ & Remark~\ref{rem:centre} \\
$W_1(f)$ & the level-one weight $\sum_i\hat f(\{i\})^2$ & Remark~\ref{rem:centre} \\
\bottomrule
\end{longtable}

\begin{thebibliography}{99}
\bibitem{FK96} E.~Friedgut, G.~Kalai, Every monotone graph property has a sharp threshold, Proc. Amer. Math. Soc.
124 (1996), 2993--3002.
\bibitem{CKLS} S.~Chakraborty, R.~Kulkarni, S.~V.~Lokam, N.~Saurabh, Upper bounds on Fourier entropy, Theoret.
Comput. Sci. 654 (2016), 92--112.
\bibitem{GMP} M.~J.~Gonz\'alez, P.~MacManus, M.~C.~Pereyra, Further evidence towards the Fourier
entropy-influence conjecture, arXiv:2606.00246, 2026.
\bibitem{Han} X.~Han, A new bound for the Fourier-entropy-influence conjecture, Combinatorica 45 (2025), no.~1,
Art.~4, doi:10.1007/s00493-024-00133-z; arXiv:2312.08271.
\bibitem{KKLMS} E.~Kelman, G.~Kindler, N.~Lifshitz, D.~Minzer, M.~Safra, Towards a proof of the
Fourier-entropy conjecture?, Geom. Funct. Anal. 30 (2020), 1097--1138.
\bibitem{LH25} P.~Li, G.~Han, Strengthening Han's Fourier entropy-influence inequality via an
information-theoretic proof, arXiv:2512.03117, 2025.
\bibitem{OD14} R.~O'Donnell, \emph{Analysis of Boolean Functions}, Cambridge University Press, 2014.
\bibitem{OT13} R.~O'Donnell, L.-Y.~Tan, A composition theorem for the Fourier entropy-influence conjecture, ICALP 2013,
Lecture Notes in Comput. Sci. 7965, 780--791, doi:10.1007/978-3-642-39206-1\_66.
\bibitem{OWZ11} R.~O'Donnell, J.~Wright, Y.~Zhou, The Fourier entropy-influence conjecture for certain classes of
Boolean functions, ICALP 2011, Lecture Notes in Comput. Sci. 6755, 330--341.
\bibitem{Shm1} Y.~Shmalo, The one-coordinate step in the Fourier entropy-influence conjecture,
\url{https://github.com/yspennstate/fourier-entropy-influence-one-coordinate}, doi:\href{https://doi.org/10.5281/zenodo.24125404}{10.5281/zenodo.24125404}, 2026.
\bibitem{Shm2} Y.~Shmalo, Fourier entropy of threshold functions and their compositions, arXiv:2610.16069, 2026.
\bibitem{VW08} S.~Verd\'u, T.~Weissman, The information lost in erasures, IEEE Trans. Inform. Theory 54 (2008),
5030--5058.
\bibitem{WWW} A.~Wan, J.~Wright, C.~Wu, Decision trees, protocols, and the Fourier entropy-influence conjecture, ITCS 2014,
arXiv:1312.3003.
\bibitem{MathBrainRepo} Y.~Shmalo, \emph{Math Brain: public snapshot, code and tests}, 2026,
\url{https://github.com/yspennstate/math-brain-public}.
\bibitem{MathBrainPaper} Y.~Shmalo, \emph{Math Brain: shared, auditable memory for multi-agent mathematical research},
October 2026, \url{https://github.com/yspennstate/math-brain-public/tree/main/papers/math-brain-2026-10-24}.
\bibitem{BrainFramework} Y.~Shmalo, \emph{Brain Framework}, public software repository, 2026,
\url{https://github.com/yspennstate/brain-framework-public}.
\end{thebibliography}

\end{document}
